\chapter{Erzeugte Ordnungen und zusätzliche Vergleiche}
\label{b12guide:ch:generated}

\section{Wann Vergleichsvorgaben eine Halbordnung erzeugen}

\begin{definition}[Halbordnungserweiterung und endliche Vergleichsfolgen]
\label{b12guide:def:generated}
Für eine beliebige Relation $E\subseteq P\times P$ sei
\[
 \mathcal E_P(E)=\{R\in\PO(P)\mid E\subseteq R\}.
\]
Ein Mitglied dieser Familie heißt \emph{Halbordnungserweiterung}
von $E$ auf demselben Träger.
Die \emph{reflexive transitive Hülle} $E^*$ wird definiert durch
\[
\begin{aligned}
 (x,y)\in E^*\quad\Longleftrightarrow\quad
 &\exists n\geq0\;\exists x_0,\ldots,x_n\in P:\\[-.3em]
 &x_0=x,\ x_n=y,\quad
   (x_i,x_{i+1})\in E\ (0\leq i<n).
\end{aligned}
\]
Hier ist $n$ eine natürliche Zahl; für $n=0$ bedeutet die
Bedingung $x=y$. Falls $\mathcal E_P(E)$ ein kleinstes Mitglied
bezüglich Inklusion besitzt, heißt es die \emph{von $E$ erzeugte
Halbordnung}.
\end{definition}

\begin{theorem}[Existenzkriterium ohne nichttriviale Kreise]
\label{b12guide:thm:extension-exists}
Für $E\subseteq P\times P$ sind äquivalent:
\begin{enumerate}
\item $\mathcal E_P(E)\ne\varnothing$.
\item $E^*$ ist antisymmetrisch.
\item Jede geschlossene endliche Vergleichsfolge
$x_0,\ldots,x_n=x_0$ mit $n\geq1$ und
$(x_i,x_{i+1})\in E$ für $0\leq i<n$ besteht nur aus gleichen
Elementen.
\end{enumerate}
Reflexive Paare $(x,x)$ sind in $E$ zugelassen; verboten sind
Kreise, die verschiedene Elemente verbinden.
\end{theorem}

\begin{theorem}[Kleinste erzeugte Halbordnung]
\label{b12guide:thm:generated}
Ist $\mathcal E_P(E)\ne\varnothing$, so ist $E^*$ die eindeutige
kleinste Halbordnung auf $P$, die $E$ enthält. Es gilt
\[
 E^*=\bigcap\mathcal E_P(E).
\]
Für jedes $R\in\PO(P)$ gilt somit
\[
 E\subseteq R\quad\Longleftrightarrow\quad E^*\subseteq R.
\]
Die Existenzbedingung kann nicht weggelassen werden.
Insbesondere ergibt $E=\varnothing$ die diskrete Ordnung
$E^*=\Delta_P$.
\end{theorem}

\begin{theorem}[Monotonie und wiederholtes Erzeugen]
\label{b12guide:thm:generated-laws}
Seien $E\subseteq F\subseteq P\times P$ und
$\mathcal E_P(F)\ne\varnothing$. Dann gilt
\[
 \mathcal E_P(E)\ne\varnothing,\qquad E^*\subseteq F^*.
\]
Für jedes $E$ mit $\mathcal E_P(E)\ne\varnothing$ gilt
\[
 E\subseteq E^*,\qquad (E^*)^*=E^*,\qquad
 (E^{\mathrm{op}})^*=(E^*)^{\mathrm{op}}.
\]
\end{theorem}

\section{Eine einzelne Vergleichsbeziehung ergänzen}

\begin{theorem}[Explizite kleinste Erweiterung um einen Vergleich]
\label{b12guide:thm:add-comparison}
Seien $R\in\PO(P)$ und $a,b\in P$ mit $a\ne b$ und
$b\nleq_R a$. Dann ist
\[
 R[a\leq b]
   =R\cup\bigl((\downarrow_R a)\times(\uparrow_R b)\bigr)
\]
die kleinste Halbordnung auf $P$, die $R$ und das Paar $(a,b)$
enthält. Explizit gilt
\[
 x\leq_{R[a\leq b]}y
 \quad\Longleftrightarrow\quad
 x\leq_R y\ \lor\ (x\leq_R a\ \land\ b\leq_R y).
\]
Gilt bereits $a\leq_R b$, so ist $R[a\leq b]=R$.
Gilt dagegen $b<_R a$, so existiert keine Halbordnung auf $P$,
die gleichzeitig $R$ und $(a,b)$ enthält.
Für $a=b$ ist keine Erweiterung nötig.
\end{theorem}

\begin{theorem}[Beide Orientierungen eines unvergleichbaren Paares]
\label{b12guide:thm:incomparable-extension}
Sind $a,b\in P$ bezüglich $R\in\PO(P)$ unvergleichbar, so
existieren die beiden Halbordnungen
$R[a\leq b]$ und $R[b\leq a]$. Es gilt
\[
 R\subsetneq R[a\leq b],\qquad
 R\subsetneq R[b\leq a],\qquad
 R[a\leq b]\cap R[b\leq a]=R.
\]
Dies betrifft jeweils einen einzelnen zusätzlichen Vergleich;
eine Erweiterung zu einer totalen Ordnung wird nicht behauptet.
\end{theorem}
